% Chapter 5: Triple Quantum Schubert Positivity Conjecture
% Based on Mihalcea (2007) \cite{Mihalcea} and Gao-Xiong (2025) \cite{GX2025}
% This chapter records the open conjecture

\chapter{Triple Quantum Schubert Positivity Conjecture}
\label{chap:TripleQuantumPositivity}

\section{Introduction: From Double to Triple}

In \cref{chap:EquivariantQuantumPositivity}, Mihalcea proved the positivity of structure constants for the product of two Schubert classes in equivariant quantum cohomology \cite{Mihalcea}:

\[
\sigma(u) \circ \sigma(v) = \sum_{d} \sum_{w} q^d c_{u,v}^{w,d} \sigma(w)
\]

where the coefficients $c_{u,v}^{w,d} \in \Lambda = \mathbb{Z}[x_1, \ldots, x_r]$ are polynomials with nonnegative coefficients.

A natural question arises: \textbf{Does the positivity persist when the product is extended to three Schubert classes?}

\section{The Triple Quantum Schubert Positivity Conjecture}

\begin{Conjecture}[Triple Quantum Schubert Positivity Conjecture]
\label{conj:TripleQuantumPositivity}
Let $X = G/P$ be a homogeneous space, and $u, v, w \in W^P$. In equivariant quantum cohomology, the product of three Schubert classes satisfies

\[
\sigma(u) \circ \sigma(v) \circ \sigma(w) = \sum_{d} \sum_{z} q^d c_{u,v,w}^{z,d} \sigma(z)
\]

where $c_{u,v,w}^{z,d} \in \Lambda = \mathbb{Z}[x_1, \ldots, x_r]$ are polynomials in the negative roots $x_i$ with nonnegative integer coefficients.
\end{Conjecture}

\textbf{Relationship with proved cases}:

\begin{enumerate}
\item \textbf{Double case} (\cite{Mihalcea}): When the product involves two factors, the positivity of $c_{u,v}^{w,d}$ has been proved.
\item \textbf{Triple classical case} (\cite{GX2025}): When $q = 0$ (no quantum correction), the triple product coefficients $c_{u,v,w}^{z}$ are nonnegative polynomials in $t_i - y_j$.
\item \textbf{Triple quantum case} (this conjecture): When the full quantum correction $q^d$ is included, the positivity of $c_{u,v,w}^{z,d}$ remains an open problem.
\end{enumerate}

\section{Geometric Meaning of the Conjecture}

From a geometric perspective, triple quantum positivity concerns the intersection counting of three Schubert varieties:

\[
c_{u,v,w}^{z,d} = \langle \sigma(u), \sigma(v), \sigma(w) \rangle_{0,3,d}
\]

where:
\begin{itemize}
\item The notation $(0,3,d)$ indicates genus 0, 3 marked points, and curve class $d$
\item It counts the number of stable maps $f: \mathbb{P}^1 \to X$ satisfying certain conditions
\item When $d = 0$, this reduces to point intersection counting in the flag variety
\end{itemize}

\textbf{Core difficulty}: In the triple case, the configuration space of curves becomes significantly more complex. Both known positivity techniques---Graham's geometric reduction and Mihalcea's projection formula---require further development to handle three factors.

\section{Current State of Research}

As of this writing, the Triple Quantum Schubert Positivity Conjecture remains an open problem. Related progress includes:

\begin{itemize}
\item Mihalcea's work on the double case proves that positivity persists under quantum corrections \cite{Mihalcea}
\item Gao-Xiong's work on the triple classical case establishes the framework of triple Schubert calculus \cite{GX2025}
\item Unifying these two results and proving the full triple quantum positivity is a frontier topic in algebraic geometry and representation theory
\end{itemize}

\section{Remarks}

The purpose of this chapter is to record this conjecture for future reference. The completion of its proof would unify the two major positivity results in Schubert calculus, holding significant theoretical importance.

\end{chapter*}
